\documentclass[a4paper,14pt]{extarticle}
\usepackage{cmap}
\usepackage[T2A]{fontenc}
\usepackage[utf8]{inputenc}
\usepackage[russian]{babel}
\usepackage{geometry}
\geometry{left=30mm, right=15mm, top=20mm, bottom=20mm}
\usepackage{enumitem}
\usepackage{indentfirst}

\begin{document}

\section*{Цель и декомпозиция на задачи}

\textbf{Цель:} Разработка эвристических алгоритмов поиска по инвертированному индексу для ускорения процесса выбора базовых чанков в системах дедупликации данных на основе сходства (SBC) без существенной потери коэффициента сжатия.

\textbf{Декомпозиция на задачи:}

\begin{enumerate}
    \item \textbf{Математическое обоснование и разработка эвристических алгоритмов поиска кандидатов}
    \begin{itemize}
        \item Разработка математически обоснованного подхода к выбору пороговых значений (например, вычисление безопасного значения для \texttt{min\_jaccard} на основе длины перекрытий и размера списков), чтобы не отсечь потенциально хорошие кандидаты.
        \item Реализация эвристики \texttt{max\_df} и статистический подбор порога для исключения признаков со слишком длинными списками инвертированного индекса (posting lists).
        \item Реализация обхода признаков по возрастанию их частоты (\texttt{rare\_first}).
        \item Добавление математических оценок для раннего отсечения кандидатов, чей максимальный overlap не позволяет достичь порога сходства (\texttt{min\_jaccard}).
        \item Разработка механизма обратной связи для фильтрации ложноположительных ключей по метрике неудачных верификаций (\texttt{fp\_metric\_threshold}).
    \end{itemize}
    
    \item \textbf{Интеграция эвристик в конвейер дедупликации}
    \begin{itemize}
        \item Внедрение поддержки оценки частоты признака (\texttt{posting\_len}) в \texttt{index\_storage.rs}.
        \item Добавление API для регистрации успешных и ложноположительных совпадений (\texttt{record\_false\_positive}, \texttt{record\_verified}).
    \end{itemize}
    
    \item \textbf{Тестирование и бенчмаркинг}
    \begin{itemize}
        \item Подготовка репрезентативных наборов данных для тестирования.
        \item Проведение сравнительных тестов: замер скорости поиска кандидатов и влияния внедренных эвристик на итоговый коэффициент дедупликации с учетом подобранных параметров.
    \end{itemize}
    
    \item \textbf{Анализ результатов}
    \begin{itemize}
        \item Оценка эффективности внедрённых оптимизаций на основе собранных метрик.
    \end{itemize}
\end{enumerate}

\end{document}
